\clearpage\pagestyle{otMain}\setlength{\otBlocEspacement}{2.5mm}\fontsize{11}{15}\selectfont\corrigesChapitre{18}{Solides, assemblages et volumes}
\begin{blocLivre}{cor-18-1}\begin{corrige}{1}{Nommer les solides}\demarche{Observer les bases et les faces.}
1 : $\res{\text{pyramide}}$ à base carrée. 2 : $\res{\text{cylindre}}$. 3 : $\res{\text{cône}}$. 4 : $\res{\text{prisme droit}}$ à base triangulaire. 5 : $\res{\text{boule}}$. 6 : $\res{\text{pavé droit}}$.
\label{fig-cor-18-1}\figureGeo[corrige]{sol-1}
\end{corrige}\end{blocLivre}
\begin{blocLivre}{cor-18-2}\begin{corrige}{2}{Faces, arêtes, sommets}\demarche{Compter de façon organisée : base, puis côtés.}
Pavé droit : $\res{6}$ faces, $\res{12}$ arêtes, $\res{8}$ sommets. Pyramide : $\res5$ faces, $\res8$ arêtes, $\res5$ sommets.
\label{fig-cor-18-2}\figureGeo[corrige]{sol-2}
\end{corrige}\end{blocLivre}
\begin{blocLivre}{cor-18-3}\begin{corrige}{3}{Un petit cube de cubes}\demarche{Compter étage par étage.}
Deux étages de $2\times2=4$ cubes : $\res8$ cubes, soit un volume de $\res{8\ \mathrm{cm}^3}$.
\label{fig-cor-18-3}\figureGeo[corrige]{sol-3}
\end{corrige}\end{blocLivre}
\begin{blocLivre}{cor-18-4}\begin{corrige}{4}{La bonne vue de dessus}\demarche{Imaginer le regard d’en haut.}
De dessus, on voit une case pour chaque colonne de cubes : trois cases en ligne et une case en arrière, à gauche. C’est la $\res{\text{vue 2}}$.
\label{fig-cor-18-4}\figureGeo[corrige]{sol-4}
\end{corrige}\end{blocLivre}
\begin{blocLivre}{cor-18-5}\begin{corrige}{5}{L’escalier}\demarche{Compter colonne par colonne.}
Les trois marches comptent $3+2+1=\res6$ cubes : volume $\res{6\ \mathrm{cm}^3}$.
\label{fig-cor-18-5}\figureGeo[corrige]{sol-5}
\end{corrige}\end{blocLivre}
\begin{blocLivre}{cor-18-6}\begin{corrige}{6}{Les cubes cachés}\demarche{Compter étage par étage.}
$9+4+1=\res{14}$ cubes. Les 4 cubes du premier étage situés sous le carré de 2 sur 2 ne se voient pas : $\res4$ cubes sont cachés, et l’on en voit 10.
\label{fig-cor-18-6}\figureGeo[corrige]{sol-6}
\end{corrige}\end{blocLivre}
\begin{blocLivre}{cor-18-7}\begin{corrige}{7}{Dessiner les vues}\demarche{Une case par colonne visible dans la direction du regard.}
Les trois vues sont en vert : de face, deux colonnes de hauteurs 2 et 1 ; de dessus, un carré de 2 sur 2 ; de droite, deux colonnes de hauteurs 2 et 1.
\label{fig-cor-18-7}\figureGeo[corrige]{sol-7}
\end{corrige}\end{blocLivre}
\begin{blocLivre}{cor-18-8}\begin{corrige}{8}{Le plan coté}\demarche{Additionner les nombres des cases.}
$3+2+1+2+1+1=\res{10}$ cubes, soit $\res{10\ \mathrm{cm}^3}$.
\label{fig-cor-18-8}\figureGeo[corrige]{sol-8}
\end{corrige}\end{blocLivre}
\begin{blocLivre}{cor-18-9}\begin{corrige}{9}{Comparer deux volumes}\demarche{Compter les deux assemblages.}
A : $2\times3+3=9$ cubes. B : $2\times2\times2=8$ cubes. $\res{\text{A}}$ a le plus grand volume, d’un cube.
\label{fig-cor-18-9}\figureGeo[corrige]{sol-9}
\end{corrige}\end{blocLivre}
\begin{blocLivre}{cor-18-10}\begin{corrige}{10}{Compléter en grand cube}\demarche{Comparer au cube complet.}
Un cube de 3 sur 3 sur 3 contient $3\times3\times3=27$ cubes : il faut en ajouter $27-20=\res7$.
\label{fig-cor-18-10}\figureGeo[corrige]{sol-10}
\end{corrige}\end{blocLivre}
\begin{blocLivre}{cor-18-11}\begin{corrige}{11}{La boîte de sucres}\demarche{Compter un étage, puis multiplier par le nombre d’étages.}
Un étage : $4\times3=12$ sucres ; deux étages : $\res{24}$ sucres. Volume : $\res{24\ \mathrm{cm}^3}$.
\label{fig-cor-18-11}\figureGeo[corrige]{sol-11}
\end{corrige}\end{blocLivre}
\begin{blocLivre}{cor-18-12}\begin{corrige}{12}{Le mur percé}\demarche{Compléter, puis retirer.}
Mur plein : $6\times4=24$ cubes. Moins la fenêtre : $24-2=\res{22}$ cubes.
\label{fig-cor-18-12}\figureGeo[corrige]{sol-12}
\end{corrige}\end{blocLivre}
\begin{blocLivre}{cor-18-13}\begin{corrige}{13}{Le grand cube peint}\demarche{Classer les petits cubes selon leur position.}
Les 8 cubes des coins ont $\res3$ faces peintes ($\res8$ cubes). Les cubes du milieu des 12 arêtes en ont 2 ($\res{12}$ cubes). Les centres des 6 faces en ont 1 ($\res6$ cubes). Le cube central n’en a aucune ($\res1$ cube). Contrôle : $8+12+6+1=27$.
\label{fig-cor-18-13}\figureGeo[corrige]{sol-13}
\end{corrige}\end{blocLivre}
\begin{blocLivre}{cor-18-14}\begin{corrige}{14}{Des vues à l’assemblage}\demarche{Croiser les informations des trois vues.}
Vue de dessus : 4 colonnes. La vue de face et la vue de droite n’autorisent qu’une seule colonne de 2 cubes, en arrière à gauche ; les trois autres colonnes ont 1 cube. Total : $2+1+1+1=\res5$ cubes.
\label{fig-cor-18-14}\figureGeo[corrige]{sol-14}
\end{corrige}\end{blocLivre}
\begin{blocLivre}{cor-18-15}\begin{corrige}{15}{Même volume ?}\demarche{Compter les cubes unités.}
Chacun compte 6 cubes : ils ont le même volume, $\res{6\ \mathrm{cm}^3}$. Elle a $\res{\text{tort}}$ : la forme change, le volume non.
\label{fig-cor-18-15}\figureGeo[corrige]{sol-15}
\end{corrige}\end{blocLivre}
